\documentclass[00-main.tex]{subfiles}

\begin{document}

\section{Validation by Simulation}\label{sec:simulations}
\subsection{Simulation Settings}\label{sec:simulation-settings}
The finite-sample performance of the pseudo-MLE $\hat{\bm\beta}^*$ \eqref{eq:beta_}, equivalently $\hat{\bm\theta}^*$ without covariates,  and the covariance estimators $\widehat{V}_\text{AR}$--$\widehat{V}_\text{OP}$ \eqref{eq:V_ar}--\eqref{eq:V_op} was assessed via Monte Carlo simulations across three settings. In each setting, $b =$ 900 independent replications of a sequence of choice and reaction times $\bm{Y} = (Y_1, \dots, Y_n)$ of length $n =$ 1000 were simulated under the true density ($f_0$; \ref{eq:f_dep}) and fit under the working density ($f^*$; \ref{eq:f_ind}). Settings vary along two dimensions: specification ($f_0 = f^*$ or $f_0 \neq f^*$) and parameter structure (fixed or varying across trials), which together determine whether the information matrix equality $\mathcal{H}_{\bm\beta} = \mathscr{I}_{\bm\beta}$ holds and which covariance estimators are consistent (Section~\ref{sec:covariance-estimation}). The first two settings are correctly specified, the second additionally demonstrates modeling of nonstationarity through covariates, while the third provides a realistic test of robustness when parameter variability is present but unmodeled.

\begin{itemize}

    \item \emph{Correct specification, constant parameters} ($f_0 = f^*$): a baseline correctly specified setting with parameters fixed across trials ($\bm{\theta}_i = \bm{\theta}^*$, no covariates). Decision outcomes are $iid$, the information matrix equality holds, and all four covariance estimators are expected to be consistent.

    \item \emph{Correct specification, trending drift rate} ($f_0 = f^*$): a correctly specified setting where the drift rate depends linearly on motion coherence via $v_i = g_v(X_i, \beta_v) = \beta_v \mathrm{coherence}_i$, with coherence decreasing monotonically with trial $i$, while all other parameters are held constant. The information matrix equality holds, and all four estimators are expected to be consistent. This setting demonstrates two capabilities of the pseudo-MLE: modeling DDM parameters as functions of covariates (Section~\ref{sec:covariates}), and modeling nonstationary parameter trends with covariates so that the joint process $(X_i, Y_i)$ remains stationary (Section~\ref{sec:assumptions}).

    \item \emph{Misspecification, unmodeled parameter variability} ($f_0 \neq f^*$): a misspecified setting where parameters $t_0,v,z$ vary $iid$ across trials, but the working model assumes constant parameters. This constitutes misspecification of type (i; Section~\ref{sec:misspecification}), where the information matrix equality fails and only the sandwich-form estimators $\widehat{V}_\text{MR}$ and $\widehat{V}_\text{AR}$ are expected to remain consistent for the covariance around the pseudo-true parameter $\bm{\theta}^*$.

\end{itemize}

All three settings use realistic parameters taken from \citet{matzke_psychological_2009} (Table~3 Means), rescaled to the present implementation: $\bar{a} = 0.63$, $\bar{t}_0 = 0.435 \pm 0.183$, $\bar{v} = 2.23 \pm 1.33$, $\bar{z} = 0.008 \pm 0.37$, where $\bar{\bm\theta}$ denotes the means of the generating distributions and $\pm$ shows variability parameters for $t_0, z$ (uniform) and $v$ (normal) under misspecification. Under correct specification, these mean values serve directly as the true parameters $\bm\theta_0$. Under misspecification, however, the pseudo-true parameters $\bm\theta^*$ are the minimizers of the divergence between the true variable-parameter density and the working constant-parameter density, and therefore need not coincide with the means $\bar{\bm\theta}$ of the generating distributions. Accordingly, $\bm\theta^*$ for the misspecification setting were determined empirically by fitting the working model to a single large sequence of $n =$ 10,000 trials; the smaller $b =$ 900 replications of $n =$ 1000 trials were then evaluated relative to this reference. Empirical standard deviations and pairwise correlations of $\hat{\bm\beta}^*$ across replications served as ground-truth standard errors and correlations for evaluating the covariance estimators.

\subsection{Simulation Results}

\subsubsection{Correct specification, constant parameters}
Simulations with constant parameters (Figure~\ref{fig:sim_constant}A,B) serve as a baseline for correct specification ($f_0 = f^*$). Point estimates (A; diagonals) are unbiased, meaning the sampling distributions are centered on the true parameters. Pairwise correlations (A; off-diagonals) are positive for $(a, v)$ and $(t_0, z)$, negative for $(a, t_0)$, $(t_0, v)$, and $(v, z)$, and near zero for $(a, z)$. Analytic standard error estimates (B; diagonals) are mostly unbiased except the non-robust estimators ($\widehat{V}_\text{SH}$, $\widehat{V}_\text{OP}$) show a slight upward bias for $v$, and the robust estimators ($\widehat{V}_\text{AR}$, $\widehat{V}_\text{MR}$) a slight downward bias for $t_0$ and $z$. The robust estimators also exhibit higher sampling variance than the non-robust estimators, most notably relative to $\widehat{V}_\text{OP}$, which has the lowest variance overall. Analytic pairwise correlation estimates (B; off-diagonals) are generally unbiased, with slight downward biases for $(a, v)$, $(t_0, z)$, and $(v, z)$ across both robust and non-robust estimators.

\begin{SCfigure}[][!h]
    \centering
    \includegraphics[width=0.5\textwidth]{docs/manuscript/figures/fig02-sim.pdf}
    \caption{\emph{Correct specification, constant parameters.}
    (A) \emph{Diagonals} show sampling distributions of $b=$ 900 point estimates ($\hat{\bm\theta}$; black density curves) around the true parameters ($\bm\theta$; blue vertical lines), where $\pm$ empirical standard deviations (red solid horizontal lines) are taken to be the true standard errors of the estimator. \emph{Off-diagonals} show bivariate sampling distributions of the point estimates (black density contours), where empirical correlations (red dashed lines) are taken to be the true correlations of the estimator. The schematic below each column indicates how the parameter varies over time (i.e., over the $n=$ 1000 trials), in this case indicating that all parameters are constant.
    (B) \emph{Diagonals} show sampling distributions of the $b=$ 900 analytic standard error estimates (i.e., square roots of the diagonal covariance elements; black=$\widehat{V}_\text{AR}$, dark gray = $\widehat{V}_\text{MR}$, gray = $\widehat{V}_\text{SH}$, light gray = $\widehat{V}_\text{OP}$) and the empirical standard errors from (A; solid red vertical lines). \emph{Off-diagonals} show analytic pairwise correlation estimates (black-gray density curves) and the empirical correlations from (A; red dashed vertical lines).}
    \label{fig:sim_constant}
\end{SCfigure}

\subsubsection{Correct specification, trending drift rate}
Simulations where the drift rate $v_i = \beta_v \cdot \mathrm{coherence}_i$ decreases monotonically with trial index (Figure~\ref{fig:sim_linear}A,B) are also correctly specified ($f_0 = f^*$) because nonstationarity is explicitly modeled by the coherence covariate. Results are mostly comparable to the constant parameter simulation above (Figure~\ref{fig:sim_constant}), with the following differences. Point estimates (A; diagonals) remain unbiased, now replacing the parameter $v$ with the coefficient $\beta_v$. Pairwise correlations (A; off-diagonals) are unchanged. Analytic standard error estimates (B; diagonals) are no longer biased upward for $\beta_v$ as the non-robust estimators were for $v$; instead, $\widehat{V}_\text{AR}$, $\widehat{V}_\text{MR}$, and $\widehat{V}_\text{OP}$ all exhibit a slight downward bias for $t_0$ and $z$. Analytic pairwise correlation estimates (B; off-diagonals) show fewer biases, with modest downward biases for $(a,t_0), (t_0, z)$. These results suggest that the covariance estimators remain mostly accurate when nonstationary trends are explicitly modeled.

\begin{SCfigure}[][h!]
    \centering
    \includegraphics[width=0.5\textwidth]{docs/manuscript/figures/fig03-sim.pdf}
    \caption{\emph{Correct specification, trending drift rate.}
    (A) As in Figure~\ref{fig:sim_constant}A, except the drift rate $v_i = \beta_v \cdot \mathrm{coherence}_i$ decreases monotonically with trial index $i$ as coherence declines (i.e., nonstationary), while the coefficient $\beta_v$ and all other parameters $a, t_0, z$ remain constant over time, such that the point estimates are for $\hat\beta_v$ instead of $\hat{v}$; as indicated by the schematic below the third column.
    (B) As in Figure~\ref{fig:sim_constant}B, except that $\beta_v$ replaces $v$.}
    \label{fig:sim_linear}
\end{SCfigure}

\subsubsection{Misspecification, unmodeled parameter variability}
Simulations under the seven-parameter DDM, in which $t_0$, $v$, and $z$ are drawn $iid$ from distributions centered on $\bar{\bm\theta}$ (Figure~\ref{fig:sim_realistic}A,B), are misspecified ($f_0 \neq f^*$) because this trial-to-trial variability is unaccounted for in the working model. This is an example of overdispersion, where there is greater variability in the decision outcomes than would be expected under the working model. Point estimates (A; diagonals) show upward bias for $t_0$ and $z$, and slight downward bias for $a$ and $v$, relative to the pseudo-true parameters $\bm\theta^*$, which do not coincide with $\bar{\bm\theta}$ ($\bar{a}=0.63$, $a^*=0.58$; $\bar{t}_0=0.435$, $t_0^*=0.435$; $\bar{v}=2.23$, $v^*=1.57$; $\bar{z}=0.008$, $z^*=0.075$). This discrepancy is most apparent for $v$ and $z$, which are negatively correlated even under correct specification (Figure~\ref{fig:sim_constant}A), such that the downward shift $v^* < \bar{v}$ induced by unmodeled variability is compensated by an upward shift $z^* > \bar{z}$. Pairwise correlations (A; off-diagonals) are mostly unchanged, except the positive association for $(a, v)$ disappears. Analytic standard error estimates (B; diagonals) show divergence between non-robust and robust estimators, where $\widehat{V}_\text{SH}$ and $\widehat{V}_\text{OP}$ exhibit large downward biases for $t_0$, $v$, and $z$, which are mitigated, but not fully corrected, by the robust estimators $\widehat{V}_\text{AR}$ and $\widehat{V}_\text{MR}$. The robust estimators still show a slight downward bias for $z$. Analytic pairwise correlation estimates (B; off-diagonals) follow the same pattern, where non-robust estimators are biased downward for $(a, t_0)$, $(a, z)$, $(t_0, z)$ and upward for $(a, v)$, $(z, v)$, with robust estimators largely correcting these biases except for a persistent upward bias in $(z, v)$.

\begin{SCfigure}[][h!]
    \centering
    \includegraphics[width=0.5\textwidth]{docs/manuscript/figures/fig04-sim.pdf}
    \caption{\emph{Misspecification, unmodeled parameter variability.}
    (A) As in Figure~\ref{fig:sim_constant}A, except the blue vertical lines now indicate pseudo-true parameters under misspecification. In this simulation, the parameter $a$ is constant across time while $t_0$, $v$, and $z$ vary $iid$ across time, as indicated by the schematic below each column.
    (B) As in Figure~\ref{fig:sim_constant}B, except that under this condition, only the black=$\widehat{V}_\text{AR}$ and dark gray=$\widehat{V}_\text{MR}$ are expected to be consistent estimators under misspecification (gray=$\widehat{V}_\text{SH}$, light gray=$\widehat{V}_\text{OP}$ are not).}
    \label{fig:sim_realistic}
\end{SCfigure}

Overall, across the three simulations, the proposed pseudo-MLE with robust covariances recovered both point estimates and standard errors with minimal bias, implying that the resulting confidence intervals achieve approximately nominal coverage under realistic temporal dependence and parameter variability. This provides a theoretically justified and empirically validated way to fit DDM parameters for which point estimates and confidence intervals remain trustworthy even in the face of unmodeled variability and misspecification, making it possible to distinguish genuine temporal structure in decision states from estimation noise. Such calibrated uncertainty quantification under realistic temporal structure is not available from standard DDM fitting approaches that assume independent, stationary trials, highlighting a capability unique to the proposed method.

\end{document}
